\chapter{Nine Counting Sequences from Horimetrik}
\label{ch:zaehlungen}

% Register: D8, D9, S9, S19, S22, S41, V11. Sources: Appendix A.4-A.6,
% C.6, F.1-F.5, H, oeis/ENTWUERFE.md.

Theories leave fingerprints: number sequences that arise when you
count their objects. In the course of this project Horimetrik has
produced nine such sequences, and not one of them appears in the
OEIS, the great encyclopaedia of integer sequences -- which is no
proof of novelty, but an indication that nobody has counted here
before. All nine are being submitted to the OEIS; once the
A-numbers have been assigned, we will enter them
here.\footnote{Placeholder: A??????. The submission drafts, with
jargon-free definitions and programs, are in the project folder
\texttt{oeis/}.}

\section{The profiles}

\paragraph{The jump positions of the depth law.}
$3,\ 15,\ 84,\ 533,\ 3883,\ 31998,\ 294875,\ 3006206,\ \ldots$
The shell in which the commutator depth
(Theorem~\ref{satz:tiefengesetz}) first reaches $-g$.
The queen of our sequences: for it an asymptotic law is
\emph{proved} ($m_g/(g+2)!\to1$), and it has passed three
prediction tests -- the last one exactly: $m_7=294875$ was
forecast before it was computed.

\paragraph{The maximal product height of squares.}
$1,\ 6,\ 22,\ 87,\ 407,\ 2473,\ 19527,\ \ldots$
How large a horizon can become when you multiply two maximally
filled shapes of capacity $r$ (register D8). Meanwhile proved: the
sequence grows like $r!\cdot e^{2\sqrt r+O(\log r)}$ -- a factorial
times a gently exploding factor (Appendix~\ref{ch:anhang-beweise});
only the fine structure of the logarithm term is still open.

\paragraph{The shell jumpers.}
$1,\ 8,\ 60,\ 360,\ 2520,\ 21168,\ 211680,\ 2268000,\ \ldots$
The number of values that a single leading zero lifts into the
next factorial shell. This sequence gave us the $n!/2$ debacle
(Chapter~\ref{ch:sackgassen}): for three terms it looked like a
half-factorial, at the fourth it no longer did. Its complement,
the \emph{non-jumpers} $n!-D(n)$
($5,\ 16,\ 60,\ 360,\ 2520,\ 19152,\ \ldots$), counts as a
sequence of the collection in its own right -- distinguished by a
fully proved closed formula, the digit formula from
Appendix~\ref{ch:anhang-beweise}.

\paragraph{The two-sided fixed points.}
$4,\ 17,\ 88,\ 540,\ 3960,\ 32760,\ \ldots$
Numbers whose representation is simultaneously stable under both
tidying-up operations (Section~\ref{sec:kommutator}); countable as
$n\cdot n!$ minus shell jumpers, for among the $n\cdot n!$
value-compact elements of a shell, exactly those without a digit
at its cap fail structure-compactness -- and those are the jumper
sequences.

\paragraph{The single-digit factorials.}
$1,\ 2,\ 5,\ 7,\ 11,\ 23,\ 29,\ 39,\ 59,\ 119,\ 143,\ \ldots$
The $n$ for which $n!$ is a single digit in its own shell --
completely characterized as $n=(i+1)!/d-1$ with $1\le d\le i$,
exactly $i$ of them per position
(Appendix~\ref{ch:anhang-beweise}); position $i=1$ supplies the
trivial initial term $n=1$. The only one of our sequences that
already came with a closed description at the moment of its
discovery -- and the supplier of nine predictions that came true
exactly.

\paragraph{The composition explosion.}
$1,\ 4,\ 23,\ 6541,\ 477\,149\,751,\ 26\,594\,371\,862\,819\,905,\ \ldots$
How high the capacity can climb when a maximally filled shape of
capacity $r$ is \emph{substituted} into another of the same kind
(Chapter~\ref{ch:strukturpolynome}). The youngest and by far the
fastest-growing sequence of the collection: substitution
multiplies the degrees and makes the coefficients explode. Thanks
to the extremal principle it is nevertheless effortless to compute
-- a single canonical substitution per term. Its asymptotics are
open.

\paragraph{The interchange maxima and zero count.}
$0,\ 1,\ 6,\ 43,\ 351,\ 2873,\ 26443,\ 270291,\ \ldots$ and
$2,\ 5,\ 13,\ 23,\ 61,\ 111,\ 274,\ 564,\ \ldots$
The two growth operations -- first prepend a zero, then relocate
value-preservingly, or the other way round -- do not commute. The
difference is measured by the \emph{interchange defect} $\Omega$:
the value after ``grow first, then lift'' minus the value after
``lift first, then grow''. This defect is \emph{never negative}:
growing first, then lifting never yields less. The smallest
instructive example is the number three as $10$ in $H_2$. Grow
first: $010$ in $H_3$ is the \emph{four}; lifted
value-preservingly it stays the four. Lift first: the three
becomes $003$; a zero in front is free, it stays the three. Hence
$\Omega=4-3=1$. This was long a conjecture and is by now a theorem
-- the proof runs through a multiplication lemma with a sharp
regime boundary. The equality cases, too, have meanwhile been
characterized: a number is indifferent to the order exactly when
both proof steps are sharp -- carry-free plus a sharp lemma -- and
the zero count obeys a semi-closed counting formula derived from
this (confirmed exactly for all measured terms).

\section{Why submit sequences?}

For two reasons, a modest one and a strategic one. The modest one:
it is a small, genuine service to mathematics -- clean
definitions, programs, a proved law. The strategic one: the OEIS
is the best novelty detector in the world. Whoever computes
$3, 15, 84, 533$ in the future, anywhere, in whatever context,
will find these entries -- and with them this book. Conversely,
the entries link back here. Should one of the sequences
nevertheless be unmasked one day as a known quantity in disguise,
that too is a result: then we finally know precisely which part of
Horimetrik is old mathematics in new clothes -- and which is not.

\begin{oberflaeche}
Nine counting questions, nine rows of answers that apparently
nobody has written down before. Perhaps because they are new.
Perhaps only because nobody ever asked. The difference between
these two possibilities is exactly what a submission settles --
and that is why we submit.
\end{oberflaeche}
